\section*{Capítulo \ref{ch:g10:measuring-the-earth} --- Medindo a Terra}

\begin{solution}{exo:g10:measuring-the-earth:1}
Um navio desaparece pelo casco primeiro: num mar plano ele apenas
encolheria, permanecendo inteiro visível. A sombra da Terra na Lua durante
um eclipse é sempre um arco de circunferência: um disco plano projetaria uma
sombra alongada quando visto de lado. Viajando para o norte, a Estrela
Polar sobe acima do horizonte: numa Terra plana a sua altura ficaria quase
a mesma.
\end{solution}

\begin{solution}{exo:g10:measuring-the-earth:2}
A \omterm{def:g10:measuring-the-earth:geodesy}{geodésia} mede a forma e o tamanho da Terra e as posições dos pontos da
sua superfície, e como eles mudam com o tempo. A \omterm{def:g10:measuring-the-earth:triangulation}{triangulação} localiza
pontos por uma cadeia de triângulos em que todos os ângulos são medidos;
só um comprimento, a base, é medido no chão.
\end{solution}

\begin{solution}{exo:g10:measuring-the-earth:3}
$C = 50 \times 5000 = 250\,000$ estádios; $250\,000 \times \qty{157.2}{m} =
\qty{39300}{km}$.
\end{solution}

\begin{solution}{exo:g10:measuring-the-earth:4}
$111/60 \approx \qty{1.85}{km}$. Quarto de meridiano: $90 \times 111 \approx
\qty{10000}{km}$.
\end{solution}

\begin{solution}{exo:g10:measuring-the-earth:5}
$a - b = 6378.137 - 6356.752 = \qty{21.385}{km}$;
$f = 21.385/6378.137 \approx 0.003353 \approx 1/298$.
\end{solution}

\begin{solution}{exo:g10:measuring-the-earth:6}
Sul: $\tan\theta_1 = 0.364$, $\theta_1 \approx \qty{20.0}{\degree}$. Norte:
$\tan\theta_2 = 0.456$, $\theta_2 \approx \qty{24.5}{\degree}$.
$\alpha = \qty{4.5}{\degree}$, logo $C = (360/4.5) \times 500 = 80 \times 500 =
\qty{40000}{km}$.
\end{solution}

\begin{solution}{exo:g10:measuring-the-earth:7}
$C = 180 - 62 - 71 = \qty{47}{\degree}$.
$AC = 10.00 \sin 71^\circ/\sin 47^\circ = 10.00 \times 0.9455/0.7314 \approx
\qty{12.93}{km}$; $BC = 10.00 \sin 62^\circ/\sin 47^\circ \approx
\qty{12.07}{km}$.
\end{solution}

\begin{solution}{exo:g10:measuring-the-earth:8}
A variação de longitude é $10 + 20 = \qty{30}{\degree}$. A
\qty{60}{\degree} N, um grau mede $111.3 \times 0.5 = \qty{55.65}{km}$, logo
a distância é $30 \times 55.65 \approx \qty{1670}{km}$. Sobre o equador:
$30 \times 111.3 \approx \qty{3340}{km}$, o dobro.
\end{solution}

\begin{solution}{exo:g10:measuring-the-earth:9}
A linha de visada toca o mar em $T$: uma tangente a uma circunferência é
perpendicular ao raio no ponto de contato, logo o ângulo em $T$ é reto.
Pitágoras: $(R+h)^2 = R^2 + d^2$, logo $d^2 = 2Rh + h^2$; $h$ é minúsculo
comparado com $R$, logo $d \approx \sqrt{2Rh} = \sqrt{2 \times
6\,371\,000 \times 2} \approx \qty{5050}{m}$, cerca de \qty{5}{km}.
\end{solution}

\begin{solution}{exo:g10:measuring-the-earth:10}
A curva é mais inclinada nas latitudes médias, por volta de
\qty{30}{\degree} a \qty{60}{\degree}, e plana perto do equador e do polo.
Numa Terra achatada, a superfície é menos curva perto dos polos; quanto
menos curva a superfície, maior a distância a percorrer para que a vertical
(e portanto a latitude) gire um grau.
\end{solution}

\begin{solution}{exo:g10:measuring-the-earth:11}
$1072.2/9.646 \approx \qty{111.15}{km}$ por grau; $90 \times 111.15
\approx \qty{10004}{km}$, \qty{2}{km} a mais que o verdadeiro quarto de
meridiano. O arco está centrado perto de \qty{46}{\degree}, onde um grau é
um pouco mais longo que o grau médio do equador ao polo
($10\,002.0/90 \approx \qty{111.13}{km}$).
\end{solution}

\begin{solution}{exo:g10:measuring-the-earth:12}
Com \qty{7.1}{\degree}: $C = (360/7.1) \times 5000 \approx 253\,500$
estádios em vez de $250\,000$: cerca de 1.4\,\% a mais (um ângulo menor
para a mesma distância significa uma Terra maior). A distância, contada em
passos por caminhantes ao longo do Nilo, era bem menos precisa que o
ângulo.
\end{solution}

\begin{solution}{exo:g10:measuring-the-earth:13}
Em $ABC$: $C = \qty{50}{\degree}$, $AC = 6.000 \sin 75^\circ/\sin 50^\circ
\approx \qty{7.566}{km}$. Em $ACD$: o ângulo em $D$ é $180 - 50 - 68 =
\qty{62}{\degree}$, e $AD = 7.566 \sin 68^\circ/\sin 62^\circ \approx
\qty{7.945}{km}$.
\end{solution}

\begin{solution}{exo:g10:measuring-the-earth:14}
$\cos\varphi = 1/2$ dá $\varphi = \qty{60}{\degree}$; $\cos\varphi = 0.1$
dá $\varphi \approx \qty{84.3}{\degree}$. Um grau de longitude na latitude
$\varphi$ é o grau do equador multiplicado por $\cos\varphi$: a razão
depende só do ângulo, não do raio.
\end{solution}

\begin{solution}{exo:g10:measuring-the-earth:15}
No equador: $111.195 - 110.574 = \qty{0.621}{km} = \qty{621}{m}$; no polo:
$111.694 - 111.195 = \qty{0.499}{km} = \qty{499}{m}$. Erros de várias
centenas de metros por grau, que se somam grau após grau, são muito maiores
que o metro: as posições convertidas a partir de latitudes e longitudes
precisam usar o elipsoide.
\end{solution}

\begin{solution}{pb:g10:measuring-the-earth:1}
\textbf{1.} $31.20 - 24.09 = \qty{7.11}{\degree}$.

\textbf{2.} $7.11/360 \approx 0.0198 = 1/50.6$, muito perto do $1/50$ de
Eratóstenes.

\textbf{3.} $C = 788 \times 360/7.11 \approx \qty{39900}{km}$.

\textbf{4.} O arco de meridiano entre os paralelos: a diferença dos ângulos
ao meio-dia é a diferença de latitude, que mede apenas a separação
norte--sul dos dois lugares; a componente leste--oeste da distância não
tem nenhum papel.

\textbf{5.} Ao meio-dia solar local, o Sol está no plano do meridiano, o
mesmo para os dois lugares; no mesmo dia ele tem a mesma posição entre as
estrelas, de modo que os dois ângulos só diferem por causa da curvatura da
Terra.

\textbf{6.} $180 - 58 - 66 = \qty{56}{\degree}$.

\textbf{7.} $AC = 12.000 \sin 66^\circ/\sin 56^\circ = 12.000 \times
0.9135/0.8290 \approx \qty{13.223}{km}$.

\textbf{8.} $BC = 12.000 \sin 58^\circ/\sin 56^\circ \approx \qty{12.275}{km}$.

\textbf{9.} O ângulo em $D$ é $180 - 47 - 75 = \qty{58}{\degree}$;
$AD = 13.223 \sin 75^\circ/\sin 58^\circ \approx \qty{15.061}{km}$.

\textbf{10.} $15.061 \cos 8^\circ \approx \qty{14.914}{km}$.

\textbf{11.} $14.914/111.13 \approx \qty{0.1342}{\degree}$, isto é,
$0.1342 \times 60 \approx 8.05$ minutos de arco.

\textbf{12.} $111.51 - 110.57 = \qty{0.94}{km}$: \qty{940}{m}.

\textbf{13.} Perto do polo, a vertical gira mais devagar com a distância: a
superfície ali é menos curva. A Terra é achatada nos polos e dilatada no
equador.

\textbf{14.} $a - b = \qty{21.385}{km}$; $f = 21.385/6378.137 \approx 1/298$.

\textbf{15.} O fio de prumo é perpendicular à superfície do elipsoide.
Numa elipse, a perpendicular à curva só passa pelo centro no equador e nos
polos; a \qty{45}{\degree} ela corta o eixo abaixo do centro.

\textbf{16.} $10\,000\,000$ metros, isto é, \qty{10000}{km}.

\textbf{17.} $10\,001\,966/10\,000\,000 \approx \qty{1.0001966}{m}$.

\textbf{18.} Curto demais: a verdadeira décima milionésima parte vale
$1.0001966$ do metro de 1799, logo o metro de 1799 é curto em cerca de
$1.97 \times 10^{-4}$ do seu comprimento, 0.02\,\%.

\textbf{19.} Curto demais. A barra é a décima milionésima parte do quarto de
meridiano calculado nos anos 1790; como a barra é mais curta que a décima
milionésima parte do verdadeiro quarto de meridiano, o quarto de meridiano
calculado era mais curto que o verdadeiro.

\textbf{20.} $1.0001966 - 1 = 0.0001966$ m: o primeiro metro era cerca de
\textbf{\qty{0.2}{mm} curto demais} --- cerca de duas folhas de papel em um
metro, uma precisão espantosa para um levantamento feito com lunetas e
ângulos através de um país inteiro.
\end{solution}
